\documentclass[11pt]{article}
\usepackage{icmlko}
\begin{document}

\title{규모가 부호를 바꾼다 --- 같은 가림이 갈래 안에서는 이롭고 수집
전체에서는 해롭다}
\author{월드모델 실험실}
\date{노트 468 · 2026년 8월 2일}
\maketitle

\begin{abstract}
노트 466 이 세운 규칙 --- \emph{모델이 그 축을 축값에 단조롭게 쓰면 두고
아니면 가린다} --- 을 마지막으로 시험했다. 노트 467 이 축 빼기로는
안 된다고 했으니 이번엔 \textbf{행}을 갈랐다. 갈래로 겹치지 않게 나눠
하위 수집 열하나.

\textbf{두 짝 다 반대로 간다.} 짝 안 상관이 KR $\mathbf{-1.000}$ ·
앱 $\mathbf{-0.900}$(p $0.037$), 짝 평균을 뺀 뒤 모아서
$\mathbf{-0.758}$(p $0.011$). 예측은 양수였다. \textbf{규칙은 죽는다.}

\textbf{그런데 표가 더 큰 것을 보여 준다.}

\begin{center}
\begin{tabular}{lrrr}
\hline
짝 & 수집 전체 & 무작위 겹 셋 & \textbf{갈래로 가름} \\
\hline
KR 만화 & $-0.0584$ & $-0.0585$ & $\mathbf{+0.0221}$ (양수 $5/6$) \\
비게임 앱 & $+0.0435$ & $+0.0421$ & $\mathbf{-0.0084}$ (양수 $2/5$) \\
\hline
\end{tabular}
\end{center}

\textbf{두 짝 모두 부호가 뒤집힌다.} 같은 모델, 같은 행, 같은 가림인데
\emph{어느 범위에서 순위를 매기느냐}가 부호를 정한다.
\end{abstract}

\begin{figure}[t]
\centering
\includegraphics[width=\textwidth]{figs/f1.pdf}
\caption{\textbf{가} --- 왼쪽에서 오른쪽으로 채점 범위를 좁힌다. 무작위
겹은 갈래 섞임을 보존해 전체와 같은 부호이고, 갈래로 가르면 뒤집힌다.
\textbf{나} --- 등록한 규칙은 양의 기울기를 예측했는데 두 짝 다 음이다.}
\end{figure}

\section{노트 465 가 왜 안 뒤집혔는지도 설명된다}

노트 465 는 각 짝을 \emph{무작위} 세 겹으로 갈랐고 여섯 겹이 모두 자기
수집 부호를 지켰다. 그때는 그것을 \emph{``갈림이 진짜 수집 단위''}의
증거로 읽었다. 지금 보면 더 정확한 읽기가 있다.

\textbf{무작위 겹은 갈래 섞임을 그대로 보존한다.} 겹 하나하나가 작은
KR 이고, 갈래 사이 순위 구조가 살아 있다. 그래서 부호가 안 변한다.
\textbf{갈래로 가르면 그 구조가 없어진다} --- 그래서 뒤집힌다.

\textbf{그러니까 \texttt{entry\_friction} 이 KR 에서 나르는 것은 갈래
\emph{사이} 정보다.} 가리면 갈래 사이 순위가 무너져 전체 점수가 내려가고,
갈래 \emph{안}에서는 그 축이 잡음이라 가리는 게 오히려 낫다. 앱은 정확히
반대다.

\section{그러면 ``집 밖 효과''는 한 수가 아니다}

이 실험실은 두 해 동안 집 밖 효과를 \textbf{수 하나}로 적어 왔다 ---
``KR $+0.0064$'', ``앱 $-0.0584$''. 노트 426 의 조건 ①(둘 이상 같은
부호)도, 노트 460 의 환율도, 노트 463 의 미결 판정도 전부 그 수 하나
위에 서 있다.

\textbf{그 수는 채점 범위에 딸린 값이었다.} 같은 변경이 갈래 안에서는
$+0.022$, 수집 전체에서는 $-0.058$ 이다. 어느 쪽이 옳은지는 \emph{쓰는
사람이 무엇을 줄 세우려는지}에 달렸다 --- 새 IP 하나의 값을 그 갈래
안에서 견줄 것인가, 수집 전체에서 견줄 것인가.

\begin{quote}
\textbf{조항} --- 집 밖 효과를 적을 때 \textbf{채점 범위를 같이 적는다.}
범위를 안 적은 수는 ``수집 전체''로 읽는다(지금까지 적은 것이 다 그것
이다). 범위가 다른 수는 \textbf{견주지 않는다} --- 노트 378 조건 ④의
``채점 집합이 다르면 잴 수 없음''이 집 밖에도 그대로 걸린다.
\end{quote}

\section{규칙은 죽었지만 죽은 자리가 값이다}

노트 466 의 쓰임 규칙은 반증됐다. 그런데 \textbf{반증의 모양이
가르쳐 준다} --- 쓰임이 강한 하위 수집(drama $-0.718$ · comedy $-0.549$)
일수록 가림이 \emph{이롭다}. 갈래 안에서 모델이 그 축을 세게 쓸수록
그 씀씀이가 잘못됐다는 뜻이다. \textbf{쓰임의 세기는 옳음의 표시가
아니라 위험의 표시다} --- 적어도 갈래 안에서는.

수집 전체에서는 반대였다(KR 쓰임 $-0.604$ 인데 가림이 해로움). 두
방향이 다르니 \textbf{쓰임 하나로는 어느 범위에서도 규칙을 못 만든다.}

\section*{남은 것}

\begin{itemize}
\item \textbf{판도 범위 문제를 갖는다.} 판은 도메인 \emph{안} 순위를
도메인마다 재서 가중평균한다 --- 즉 ``갈래로 가른'' 쪽에 가깝다. 집 밖
수치를 판과 견줄 때 이 차이를 안 적고 있었다.
\item \textbf{어느 범위가 배포에 맞나} --- 팝업 기획이 IP 를 고를 때
갈래 안에서 고르는지 전체에서 고르는지. 이건 자료가 아니라 쓰임새 물음
이라 사람에게 물어야 한다.
\item \texttt{OOC\_DROP} 은 여전히 미결이고, 이제 \textbf{범위마다 답이
다르다}는 것까지 안다.
\end{itemize}

\end{document}
